Требуется подобрать такую функцию $z(x)$, которая
наилучшим образом приближает экспериментальные данные.

На первом этапе производится предварительный выбор вида аппроксимирующей
функции. Для этого используются средние значения аргумента и функции:
\[
    x_a=\frac{x_1+x_9}{2}, \qquad
    x_g=\sqrt{x_1x_9}, \qquad
    x_h=\frac{2}{\frac{1}{x_1}+\frac{1}{x_9}},
\]
\[
    y_a=\frac{y_1+y_9}{2}, \qquad
    y_g=\sqrt{y_1y_9}, \qquad
    y_h=\frac{2}{\frac{1}{y_1}+\frac{1}{y_9}}.
\]

Подставим значения крайних точек:
\[
    x_1=1,\qquad x_9=5,\qquad y_1=1{,}24,\qquad y_9=7{,}07.
\]

Тогда
\[
    x_a=\frac{1+5}{2}=3,
\]
\[
    x_g=\sqrt{1\cdot 5}=\sqrt{5}\approx 2{,}2361,
\]
\[
    x_h=\frac{2}{1+\frac{1}{5}}=\frac{2}{\frac{6}{5}}=\frac{5}{3}\approx 1{,}6667.
\]

Для значений функции получаем:
\[
    y_a=\frac{1{,}24+7{,}07}{2}=\frac{8{,}31}{2}=4{,}155,
\]
\[
    y_g=\sqrt{1{,}24\cdot 7{,}07}=\sqrt{8{,}7668}\approx 2{,}9609,
\]
\[
    y_h=\frac{2}{\frac{1}{1{,}24}+\frac{1}{7{,}07}}\approx 2{,}1099.
\]

\textbf{Определение значений $z(x_a)$, $z(x_g)$, $z(x_h)$ по графику}

По построенному графику табличной функции определим значения
\[
    z(x_a),\qquad z(x_g),\qquad z(x_h).
\]

Для этого на графике отметим точки
\[
    x_a=3,\qquad x_g=\sqrt{5}\approx 2{,}2361,\qquad x_h=\frac{5}{3}\approx 1{,}6667
\]

\image{1}{График табличной функции}{1}

и проведём через них вертикальные линии до пересечения с графиком функции.
После этого из точек пересечения проведём горизонтальные линии к оси $z(x)$.

Так как график строится по табличным значениям с соединением соседних узлов
отрезками прямых, то графическое определение значений эквивалентно линейной
интерполяции.

Если точка $x^\ast$ лежит между узлами $x_i$ и $x_{i+1}$, то
\[
    z(x^\ast)=y_i+\frac{y_{i+1}-y_i}{x_{i+1}-x_i}(x^\ast-x_i).
\]

\textbf{1. Значение в точке $x_a$:}

Так как
\[
    x_a=3
\]
совпадает с табличным узлом, получаем сразу
\[
    z(x_a)=z(3)=3{,}06.
\]

\textbf{2. Значение в точке $x_g$:}

Так как
\[
    x_g\approx 2{,}2361
\]
лежит между точками $x=2$ и $x=2{,}5$, имеем
\[
    z(x_g)=1{,}61+\frac{2{,}16-1{,}61}{2{,}5-2}\,(2{,}2361-2).
\]
Отсюда
\[
    z(x_g)=1{,}61+\frac{0{,}55}{0{,}5}\cdot 0{,}2361
    =1{,}61+1{,}1\cdot 0{,}2361
    \approx 1{,}8697.
\]

\textbf{3. Значение в точке $x_h$:}

Так как
\[
    x_h\approx 1{,}6667
\]
лежит между точками $x=1{,}5$ и $x=2$, получаем
\[
    z(x_h)=1{,}74+\frac{1{,}61-1{,}74}{2-1{,}5}\,(1{,}6667-1{,}5).
\]
Следовательно,
\[
    z(x_h)=1{,}74+\frac{-0{,}13}{0{,}5}\cdot 0{,}1667
    =1{,}74-0{,}26\cdot 0{,}1667
    \approx 1{,}6967.
\]

Таким образом, по графику получаем:
\[
    z(x_a)=3{,}06,\qquad z(x_g)\approx 1{,}8697,\qquad z(x_h)\approx 1{,}6967.
\]


\textbf{Вычисление величин $\delta_1,\dots,\delta_9$}

Используем найденные ранее значения:
\[
    z(x_a)=3{,}06,\qquad z(x_g)\approx 1{,}8697,\qquad z(x_h)\approx 1{,}6967,
\]
\[
    y_a=4{,}155,\qquad y_g\approx 2{,}9609,\qquad y_h\approx 2{,}1099.
\]

Тогда получаем:

\[
    \delta_1=|z(x_a)-y_a|=|3{,}06-4{,}155|=1{,}095,
\]

\[
    \delta_2=|z(x_g)-y_g|=|1{,}8697-2{,}9609|\approx 1{,}0912,
\]

\[
    \delta_3=|z(x_a)-y_g|=|3{,}06-2{,}9609|\approx 0{,}0991,
\]

\[
    \delta_4=|z(x_g)-y_a|=|1{,}8697-4{,}155|\approx 2{,}2853,
\]

\[
    \delta_5=|z(x_h)-y_a|=|1{,}6967-4{,}155|\approx 2{,}4583,
\]

\[
    \delta_6=|z(x_a)-y_h|=|3{,}06-2{,}1099|\approx 0{,}9501,
\]

\[
    \delta_7=|z(x_h)-y_h|=|1{,}6967-2{,}1099|\approx 0{,}4133,
\]

\[
    \delta_8=|z(x_h)-y_g|=|1{,}6967-2{,}9609|\approx 1{,}2642,
\]

\[
    \delta_9=|z(x_g)-y_h|=|1{,}8697-2{,}1099|\approx 0{,}2403.
\]

Итоговая таблица
\[
    \begin{array}{c|c|c|c|c|c|c|c|c|c}
        k & 1 & 2 & 3 & 4 & 5 & 6 & 7 & 8 & 9 \\ \hline
        \delta_k & 1{,}0950 & 1{,}0912 & 0{,}0991 & 2{,}2853 & 2{,}4583 & 0{,}9501 & 0{,}4133 & 1{,}2642 & 0{,}2403
    \end{array}
\]


Минимальное значение имеет величина
\[
    \delta_3\approx 0{,}0991.
\]

Следовательно, выбираем модель третьего типа:
\[
    z_3(x)=ae^{bx}.
\]


\textbf{Определение коэффициентов модели методом наименьших квадратов}

Для выбранной модели
\[
z(x)=\alpha e^{\beta x}
\]
выполним линеаризацию:
\[
\ln z(x)=\ln \alpha+\beta x.
\]

Введём обозначения
\[
X=x,\qquad Y=\ln y.
\]
Тогда получаем линейную зависимость
\[
Y=AX+B,
\]
где
\[
A=\beta,\qquad B=\ln \alpha.
\]

Следовательно, задача сводится к нахождению коэффициентов $A$ и $B$
для линейной функции методом наименьших квадратов, то есть к минимизации
функционала
\[
\sum_{i=1}^{n}\left(AX_i+B-Y_i\right)^2 \to \min.
\]

Нормальные уравнения имеют вид
\[
A\sum_{i=1}^{n}X_i^2+B\sum_{i=1}^{n}X_i=\sum_{i=1}^{n}X_iY_i,
\]
\[
A\sum_{i=1}^{n}X_i+nB=\sum_{i=1}^{n}Y_i.
\]

Обозначим
\[
S_x=\sum_{i=1}^{n}X_i,\qquad
S_y=\sum_{i=1}^{n}Y_i,\qquad
S_{xx}=\sum_{i=1}^{n}X_i^2,\qquad
S_{xy}=\sum_{i=1}^{n}X_iY_i.
\]

Тогда коэффициенты находятся по формулам
\[
A=\frac{nS_{xy}-S_xS_y}{nS_{xx}-S_x^2},
\qquad
B=\frac{S_y-AS_x}{n}.
\]

После этого выполняем обратное преобразование:
\[
\beta=A,\qquad \alpha=e^B.
\]

Искомая аппроксимирующая функция принимает вид
\[
z(x)=\alpha e^{\beta x}.
\]

После нахождения коэффициентов вычисляется среднеквадратичное отклонение
уже для исходной функции:
\[
\Delta=\sqrt{\frac{1}{n}\sum_{i=1}^{n}\left(y_i-z(x_i)\right)^2 }.
\]